\documentclass[11pt,letterpaper]{article}

% Packages
\usepackage[margin=1in]{geometry}
\usepackage{graphicx}
\usepackage{hyperref}
\usepackage{listings}
\usepackage{xcolor}
\usepackage{fancyhdr}
\usepackage{tcolorbox}
\usepackage{enumitem}

% Color definitions
\definecolor{codebackground}{RGB}{245,245,245}
\definecolor{keywordcolor}{RGB}{0,0,255}
\definecolor{stringcolor}{RGB}{163,21,21}

% Code listing style
\lstset{
    backgroundcolor=\color{codebackground},
    basicstyle=\ttfamily\small,
    keywordstyle=\color{keywordcolor},
    stringstyle=\color{stringcolor},
    breaklines=true,
    frame=single,
    numbers=left,
    numberstyle=\tiny,
    tabsize=2
}

% Header/Footer
\pagestyle{fancy}
\fancyhf{}
\fancyhead[L]{Specs to Machine | Fusion Manufacturing Pipeline}
\fancyhead[R]{User Guide}
\fancyfoot[C]{\thepage}

% Title information
\title{\textbf{Specs to Machine | Fusion Manufacturing Pipeline}\\
\large User Guide}
\author{Bobrick Washroom Equipment\\James Derrod}
\date{\today}

\begin{document}

\maketitle
\thispagestyle{empty}

\vspace{1cm}

\newpage
\tableofcontents
\newpage

% ============================================================================
\section{Overview}

The Specs to Machine manufacturing pipeline is a fully automated system that processes manufacturing orders for bathroom partition components (doors, panels, and stiles). Orders from iBob are automatically sent to the \texttt{order\_dropbox} folder as JSON files, where the Fusion 360 extension processes them using a queue system. The system automatically generates:

\begin{itemize}
    \item 3D STEP models
    \item CNC G-code programs (.nc files)
    \item Parameter documentation (CSV and JSON)
    \item Processing logs
\end{itemize}

\subsection{Key Features}

\begin{itemize}
    \item \textbf{Fully Automated:} No manual intervention required - iBob orders are processed automatically
    \item \textbf{Queue System:} Handles multiple orders in sequence, ensuring reliable processing
    \item \textbf{Batch Processing:} Processes multiple components per order (doors, panels, stiles)
    \item \textbf{Network Storage:} All files stored on shared network drive
    \item \textbf{Comprehensive Logging:} Full audit trail for each order
    \item \textbf{Error Handling:} Failed orders moved to separate folder for review
\end{itemize}

\subsection{System Architecture}

\begin{tcolorbox}[colback=blue!5!white,colframe=blue!75!black,title=Processing Flow]
\begin{enumerate}[noitemsep]
    \item iBob generates JSON order file and places in \texttt{order\_dropbox} folder
    \item Extension detects file and adds to processing queue
    \item System processes order: apply parameters → export model → generate G-code
    \item Outputs saved to \texttt{output} folder organized by order ID
    \item Completed order moved to \texttt{order\_completed} or \texttt{order\_failed}
\end{enumerate}
\end{tcolorbox}

% ============================================================================
\section{System Requirements}

\subsection{Software Requirements}

\begin{itemize}
    \item \textbf{Autodesk Fusion 360 Desktop Edition:} Current version with CAM capabilities
    \item \textbf{Virtual Machine:} Windows 10/11 VM
    \item \textbf{Network Access:} VM must have read/write access to \texttt{M:\textbackslash S2S File Test}
\end{itemize}

\subsection{Hardware Requirements}

\begin{itemize}
    \item \textbf{Processor:} Intel Core i5 or better (i7 recommended)
    \item \textbf{RAM:} 8 GB minimum (16 GB recommended)
    \item \textbf{Storage:} Network access with adequate bandwidth
\end{itemize}

% ============================================================================
\section{Installation \& Setup (IT Department)}

\subsection{System Setup}

The Specs to Machine manufacturing pipeline runs in a Windows virtual machine with Fusion 360 Desktop Edition installed. Once configured, the system automatically monitors the network folders and processes incoming orders.

\subsection{Network Folder Structure}

The following folder structure will be created automatically on the network drive at \texttt{M:\textbackslash S2S File Test}:

\begin{lstlisting}[language=bash]
M:\S2S File Test\
├── input\
│   ├── door\
│   ├── panel\
│   └── stile\
├── order_dropbox\
├── order_processing\
├── order_completed\
├── order_failed\
└── output\
    ├── models\
    ├── gcode\
    ├── parameters\
    └── logs\
\end{lstlisting}

\subsection{Installing the Fusion 360 Extension}

\begin{enumerate}
    \item Install Fusion 360 Desktop Edition on the virtual machine
    
    \item Copy the \texttt{FusionExtension} folder to:\\
    \texttt{C:\textbackslash Users\textbackslash [username]\textbackslash AppData\textbackslash Roaming\textbackslash Autodesk\textbackslash}\\
    \texttt{Autodesk Fusion 360\textbackslash API\textbackslash AddIns\textbackslash}
    
    \item Open Fusion 360 on the virtual machine
    
    \item Click \textbf{Utilities} → \textbf{ADD-INS}
    
    \item Find \textbf{FusionManufacturingPipeline} in the list
    
    \item Click \textbf{Run} to start the extension
    
    \item Check \textbf{Run on Startup} for automatic loading on VM startup
\end{enumerate}

\begin{tcolorbox}[colback=green!5!white,colframe=green!75!black,title=Success]
If installation is successful, the extension will create all necessary network folders and display a confirmation message. The system will immediately begin monitoring the \texttt{order\_dropbox} folder for incoming orders from iBob.
\end{tcolorbox}

\begin{tcolorbox}[colback=blue!5!white,colframe=blue!75!black,title=Automatic Processing]
Once running, the extension uses a queue system to automatically process any JSON files that appear in the \texttt{order\_dropbox} folder. No manual intervention is required.
\end{tcolorbox}

\subsection{Placing Model Files}

Place the master Fusion 360 model files (.f3d) in the input folders:

\begin{itemize}
    \item \texttt{M:\textbackslash S2S File Test\textbackslash input\textbackslash door\textbackslash door.f3d}
    \item \texttt{M:\textbackslash S2S File Test\textbackslash input\textbackslash panel\textbackslash panel.f3d}
    \item \texttt{M:\textbackslash S2S File Test\textbackslash input\textbackslash stile\textbackslash stile.f3d}
\end{itemize}

\begin{tcolorbox}[colback=yellow!5!white,colframe=orange!75!black,title=Important]
Each folder should contain exactly ONE .f3d file. Multiple files may cause conflicts.
\end{tcolorbox}

\subsection{Machine Configuration}

The system is configured to use the Anderson Stratos CNC router for all G-code generation:

\begin{itemize}
    \item \textbf{Machine:} Anderson Stratos Pro Full line
    \item \textbf{Post Processor:} Anderson Stratos 2.cps (custom)
    \item \textbf{Travel:} X: 3149mm, Y: 2184mm, Z: 254mm
    \item \textbf{Spindle:} 24,000 RPM max
    \item \textbf{Tool Changer:} 10 tools
\end{itemize}

\begin{tcolorbox}[colback=blue!5!white,colframe=blue!75!black,title=CAM Setup Requirement]
The master model files (.f3d) must have CAM setups configured to use the Anderson Stratos machine definition. The extension automatically applies the Anderson Stratos 2 post processor when generating G-code - no manual post processor selection is needed.
\end{tcolorbox}

% ============================================================================
\section{How to Use (ME Department \& Business Unit)}

\subsection{Starting the System}

The system automatically starts when the extension loads on VM startup. To manually start or restart:

\begin{enumerate}
    \item Open Fusion 360 on the virtual machine
    
    \item Click the \textbf{Utilities} tab
    
    \item Locate the extension icon in the Utilities panel
    
    \item Click the extension icon to start monitoring
    
    \item Choose \textbf{YES} to start Folder Monitor mode
    
    \item The system will display: \textit{"Folder monitor started. Watching for orders..."}
\end{enumerate}

\begin{tcolorbox}[colback=blue!5!white,colframe=blue!75!black,title=Automatic Operation]
Once started, the extension continuously monitors the \texttt{order\_dropbox} folder. Files from iBob are automatically queued and processed in order. The VM should remain running 24/7 for continuous operation.
\end{tcolorbox}

\subsection{Processing an Order}

Orders are automatically submitted from iBob to the \texttt{order\_dropbox} folder. The processing is completely automated:

\begin{enumerate}
    \item iBob generates a JSON order file (see Section \ref{sec:json_format})
    
    \item iBob places the file in:\\
    \texttt{M:\textbackslash S2S File Test\textbackslash order\_dropbox\textbackslash}
    
    \item The extension automatically:
    \begin{itemize}
        \item Detects the new file using the queue system
        \item Moves it to \texttt{order\_processing}
        \item Processes each component (doors, panels, stiles)
        \item Generates all outputs (models, G-code, parameters)
        \item Moves completed order to \texttt{order\_completed} or \texttt{order\_failed}
    \end{itemize}
\end{enumerate}

\begin{tcolorbox}[colback=yellow!5!white,colframe=orange!75!black,title=Note]
No manual intervention is required. The system processes all orders from iBob automatically as they arrive in the \texttt{order\_dropbox} folder.
\end{tcolorbox}

\subsection{Finding Your Outputs}

All outputs are organized by order ID (IBUSXXXX) in the \texttt{output} folder:

\begin{lstlisting}[language=bash]
M:\S2S File Test\output\
├── models\IBUS366574\
│   ├── D1.step
│   ├── D2.step
│   └── ...
├── gcode\IBUS366574\
│   ├── 1-D1-IBUS366574.nc
│   ├── 1-D2-IBUS366574.nc
│   └── ...
├── parameters\IBUS366574\
│   ├── D1_all_parameters.csv
│   ├── D1_all_parameters.json
│   └── ...
└── logs\
    └── IBUS366574_log.txt
\end{lstlisting}

\subsection{Understanding File Names}

\begin{itemize}
    \item \textbf{Models:} \texttt{[ComponentID].step}\\
    Example: \texttt{D1.step}, \texttt{P2.step}
    
    \item \textbf{G-Code:} \texttt{1-[ComponentID]-[OrderID].nc}\\
    Example: \texttt{1-D1-IBUS366574.nc}
    
    \item \textbf{Parameters:} \texttt{[ComponentID]\_all\_parameters.[csv|json]}\\
    Example: \texttt{D1\_all\_parameters.csv}
    
    \item \textbf{Logs:} \texttt{[OrderID]\_log.txt}\\
    Example: \texttt{IBUS366574\_log.txt}
\end{itemize}

% ============================================================================
\section{JSON Order Format}
\label{sec:json_format}

\subsection{Structure Overview}

Each order file contains:
\begin{itemize}
    \item Order ID (IBUSXXXX format)
    \item List of doors (if any)
    \item List of panels (if any)
    \item List of stiles (if any)
\end{itemize}

Each parameter uses the format: \texttt{[value, datatype, description]}

\subsection{Example Order File}

\begin{lstlisting}[language=json,caption={Example Order: IBUS366574.json}]
{
  "order_id": ["IBUS366574", "string", "Order identifier"],
  "doors": [
    {
      "id": ["D1", "string", "Door ID"],
      "parameters": {
        "series_id": ["3082G.67P", "string", "Series ID"],
        "component_height": [96.0, "float", "Height in inches"],
        "component_width": [10.5, "float", "Width in inches"],
        "door_hinge_height": [75.0, "float", "Hinge height"],
        "door_keep_height": [75.0, "float", "Keep height"],
        "door_floor_clearance": [1.0, "float", "Floor clearance"]
      }
    }
  ],
  "panels": [
    {
      "id": ["P1", "string", "Panel ID"],
      "parameters": {
        "series_id": ["3082G.67P", "string", "Series ID"],
        "component_height": [97.75, "float", "Height in inches"],
        "component_width": [36.0, "float", "Width in inches"]
      }
    }
  ],
  "stiles": []
}
\end{lstlisting}

\subsection{Required Parameters by Component Type}

\subsubsection{Doors}

\begin{itemize}[noitemsep]
    \item \texttt{series\_id}: Series identifier (string)
    \item \texttt{component\_height}: Height in inches (float)
    \item \texttt{component\_width}: Width in inches (float)
    \item \texttt{door\_hinge\_height}: Hinge pilaster height (float)
    \item \texttt{door\_keep\_height}: Keep pilaster height (float)
    \item \texttt{door\_floor\_clearance}: Floor clearance (float)
\end{itemize}

\subsubsection{Panels}

\begin{itemize}[noitemsep]
    \item \texttt{series\_id}: Series identifier (string)
    \item \texttt{component\_height}: Height in inches (float)
    \item \texttt{component\_width}: Width in inches (float)
\end{itemize}

\subsubsection{Stiles}

\begin{itemize}[noitemsep]
    \item \texttt{series\_id}: Series identifier (string)
    \item \texttt{component\_height}: Height in inches (float)
    \item \texttt{component\_width}: Width in inches (float)
    \item \texttt{stile\_left\_side\_door}: Door on left side (bool)
    \item \texttt{stile\_left\_side\_hinging}: Left side is hinge stile (bool)
    \item \texttt{stile\_left\_side\_door\_height}: Left door height (float or null)
    \item \texttt{stile\_left\_side\_door\_floor\_clearance}: Left door clearance (float or null)
    \item \texttt{stile\_left\_side\_door\_swinging\_out}: Left door swings out (bool or null)
    \item \texttt{stile\_right\_side\_door}: Door on right side (bool)
    \item \texttt{stile\_right\_side\_hinging}: Right side is hinge stile (bool)
    \item \texttt{stile\_right\_side\_door\_height}: Right door height (float or null)
    \item \texttt{stile\_right\_side\_door\_floor\_clearance}: Right door clearance (float or null)
    \item \texttt{stile\_right\_side\_door\_swinging\_out}: Right door swings out (bool or null)
\end{itemize}

% ============================================================================
\section{Monitoring \& Troubleshooting}

\subsection{Monitoring Active Processing}

\begin{enumerate}
    \item Check the \texttt{order\_processing} folder - files currently being processed
    \item Check Fusion 360 - progress dialog shows current component
    \item Check \texttt{order\_completed} - successfully processed orders
    \item Check \texttt{order\_failed} - failed orders with error logs
\end{enumerate}

\subsection{Understanding Log Files}

Each order generates a detailed log file in:\\
\texttt{M:\textbackslash S2S File Test\textbackslash output\textbackslash logs\textbackslash [OrderID]\_log.txt}

The log contains:
\begin{itemize}
    \item Timestamp for each operation
    \item Parameter values being applied
    \item CAM regeneration status
    \item G-code generation results
    \item Any errors or warnings
\end{itemize}

\subsection{Common Issues}

\begin{tcolorbox}[colback=red!5!white,colframe=red!75!black,title=Problem: Order moved to failed folder]
\textbf{Solution:}
\begin{enumerate}[noitemsep]
    \item Open the corresponding \texttt{.txt} file in \texttt{order\_failed} folder
    \item Read the error message
    \item Common causes:
    \begin{itemize}
        \item Invalid JSON format
        \item Missing required parameters
        \item Model file not found in input folder
        \item CAM setup issues in master model
    \end{itemize}
    \item Fix the issue and re-submit the JSON file
\end{enumerate}
\end{tcolorbox}

\begin{tcolorbox}[colback=red!5!white,colframe=red!75!black,title=Problem: No G-code files generated]
\textbf{Solution:}
\begin{enumerate}[noitemsep]
    \item Check that CAM setups exist in the master model
    \item Verify CAM toolpaths are generated
    \item Check log file for post-processing errors
    \item Verify richauto.cps post-processor is installed in Fusion 360
\end{enumerate}
\end{tcolorbox}

\begin{tcolorbox}[colback=red!5!white,colframe=red!75!black,title=Problem: Parameters not updating correctly]
\textbf{Solution:}
\begin{enumerate}[noitemsep]
    \item Verify parameter names in JSON match model parameter names
    \item Check that parameter values are within valid ranges
    \item Review parameter CSV/JSON output to see actual values applied
    \item For stiles, ensure you're using correct stile-specific parameter names
\end{enumerate}
\end{tcolorbox}

\subsection{Restarting the System}

If the system becomes unresponsive:

\begin{enumerate}
    \item In Fusion 360, click \textbf{Utilities} → \textbf{ADD-INS}
    \item Select \textbf{FusionManufacturingPipeline}
    \item Click \textbf{Stop}
    \item Wait 3 seconds
    \item Click \textbf{Run}
    \item Start the folder monitor again
\end{enumerate}

% ============================================================================
\section{Appendix}

\subsection{Parameter Unit Conversion}

The system automatically converts all parameter values to display units:

\begin{center}
\begin{tabular}{|l|l|l|}
\hline
\textbf{Input Unit} & \textbf{Fusion Internal} & \textbf{Output Unit} \\
\hline
inches (in) & centimeters & inches (in) \\
millimeters (mm) & centimeters & millimeters (mm) \\
feet (ft) & centimeters & feet (ft) \\
\hline
\end{tabular}
\end{center}

\subsection{File Size Guidelines}

Typical file sizes:

\begin{itemize}
    \item STEP models: 500 KB - 5 MB per file
    \item G-code files: 50 KB - 500 KB per file
    \item Parameter files: 5 KB - 20 KB per file
    \item Log files: 10 KB - 100 KB per order
\end{itemize}

\subsection{Processing Time Estimates}

Per component:

\begin{itemize}
    \item Parameter application: 5-10 seconds
    \item CAM regeneration: 30-60 seconds
    \item G-code generation: 10-20 seconds
    \item Model export: 5-10 seconds
\end{itemize}

\textbf{Total per component:} Approximately 1-2 minutes

\textbf{Full order estimate:} Number of components × 1.5 minutes

% ============================================================================

\end{document}
